\chapter{Dry Orgasm}

\section*{Definition}
Orgasm achieved without ejaculation of semen (anejaculation) or with significantly reduced or absent seminal fluid, occurring retrograde into the bladder or due to failure of seminal emission.

\section*{Pathophysiology}
Dry orgasm results from either retrograde ejaculation (semen flows backward into the bladder due to failure of the bladder neck to close during emission) or anejaculation (failure of seminal emission from the seminal vesicles, prostate, and vas deferens). The internal bladder sphincter normally contracts during orgasm to prevent retrograde flow; disruption by surgery, medication, or neuropathy causes retrograde ejaculation.

\section*{Causes}
\begin{itemize}
  \item Transurethral resection of the prostate (TURP)
  \item Spinal cord injury or neurological disease
  \item Diabetes mellitus with autonomic neuropathy
  \item Medication side effects (alpha-blockers, antipsychotics, antidepressants)
  \item Retroperitoneal lymph node dissection (testicular cancer surgery)
  \item Multiple sclerosis
  \item Spinal stenosis or cauda equina syndrome
  \item Bladder neck surgery
  \item Severe hypospadias or congenital abnormalities
  \item Psychogenic anejaculation
\end{itemize}

\section*{When to See a Doctor}
See your doctor if:
\begin{itemize}
  \item Severe or worsening symptoms
  \item Persistent dry orgasm with fertility concerns, associated with back pain or neurological symptoms, or after pelvic surgery
  \item Accompanied by other concerning signs
\end{itemize}

\section*{Related Symptoms}
\begin{itemize}
  \item Erectile dysfunction
  \item Infertility
  \item Reduced libido
  \item Testicular pain
  \item Urinary symptoms
\end{itemize}

\textbf{Important:} If this symptom persists, your doctor will check for serious underlying causes.

\section*{Self-Care / Home Management}
\begin{itemize}
  \item Rest and avoid activities that make it worse.
  \item Maintain adequate hydration and a balanced diet.
  \item Over-the-counter remedies may provide symptomatic relief (consult a pharmacist or doctor).
  \item Monitor symptoms and seek professional advice if they persist or worsen.
  \item Avoid self-medicating with prescription medications without medical guidance.
\end{itemize}

\section*{How Your Doctor Will Evaluate You}
\begin{itemize}
  \item A thorough history and physical examination are the first steps in evaluation.
  \item Laboratory tests (blood work, urinalysis) may be ordered based on clinical suspicion.
  \item Imaging studies (X-ray, ultrasound, CT, MRI) may be indicated when structural pathology is suspected.
  \item Specialist referral is arranged when the diagnosis remains uncertain or requires specific expertise.
\end{itemize}

\section*{Treatment Options}
\begin{itemize}
  \item Treatment targets the underlying cause identified through diagnostic evaluation.
  \item Supportive care and symptom management are provided while awaiting definitive therapy.
  \item Medications, lifestyle modifications, and physical therapies may be prescribed as appropriate.
  \item Follow-up ensures treatment effectiveness and allows timely adjustments to the care plan.
\end{itemize}

\clearpage